\appendix
\renewcommand{\chaptername}{Anexo}
\titleformat{\chapter}[display]{\normalfont\bfseries\color{marino}\centering}{\large ANEXO \thechapter}{4pt}{\LARGE}
\lstdefinestyle{anexo}{basicstyle=\ttfamily\footnotesize\linespread{1}\selectfont,numbersep=6pt}

\chapter{Salida de \texttt{cargo test} y \texttt{cargo clippy}}\label{anx:tests}
Este anexo reproduce la salida real de las herramientas de verificación ejecutadas sobre el código del repositorio durante la compilación de este documento. El primer listado corresponde a \texttt{cargo test}, con el nombre de cada una de las 52 pruebas (51 de la biblioteca y 1 del binario), y el segundo a \texttt{cargo clippy --all-targets}, cuya ejecución concluye sin advertencias.

{\singlespacing
\lstinputlisting[language={},basicstyle=\ttfamily\scriptsize,numbers=none,caption={Salida completa de \texttt{cargo test}},label=cod:tests]{generado/cargo_test.txt}
\lstinputlisting[language={},basicstyle=\ttfamily\scriptsize,numbers=none,caption={Salida de \texttt{cargo clippy --all-targets}},label=cod:clippy]{generado/cargo_clippy.txt}
}

\chapter{Código fuente completo}\label{anx:codigo}
Este anexo contiene el código fuente completo del motor, documentado con comentarios \texttt{///} en español, junto con el archivo de configuración \texttt{Cargo.toml}. Los listados se generan automáticamente desde el repositorio, por lo que coinciden exactamente con la versión compilada y probada. En total se incluyen \LineasFuente{} líneas. El repositorio está disponible en \url{\Repositorio}.

{\singlespacing% Generado por tesis/scripts/anexo_codigo.py — no editar a mano.
\section{Configuración del paquete y biblioteca}
Esta sección contiene \texttt{Cargo.toml} (6 líneas), \texttt{src/lib.rs} (19 líneas).
\lstinputlisting[language=TOML,style=anexo,caption={Archivo \texttt{Cargo.toml}},label=cod:src:Cargo.toml]{generado/fuente/Cargo.toml}
\lstinputlisting[language=Rust,style=anexo,caption={Archivo \texttt{src/lib.rs}},label=cod:src:src:lib.rs]{generado/fuente/src/lib.rs}
\section{Módulo \texttt{math}: álgebra lineal y utilidades}
Esta sección contiene \texttt{src/math/mod.rs} (15 líneas), \texttt{src/math/error.rs} (68 líneas), \texttt{src/math/vector.rs} (176 líneas), \texttt{src/math/matrix.rs} (291 líneas), \texttt{src/math/rng.rs} (58 líneas).
\lstinputlisting[language=Rust,style=anexo,caption={Archivo \texttt{src/math/mod.rs}},label=cod:src:src:math:mod.rs]{generado/fuente/src/math/mod.rs}
\lstinputlisting[language=Rust,style=anexo,caption={Archivo \texttt{src/math/error.rs}},label=cod:src:src:math:error.rs]{generado/fuente/src/math/error.rs}
\lstinputlisting[language=Rust,style=anexo,caption={Archivo \texttt{src/math/vector.rs}},label=cod:src:src:math:vector.rs]{generado/fuente/src/math/vector.rs}
\lstinputlisting[language=Rust,style=anexo,caption={Archivo \texttt{src/math/matrix.rs}},label=cod:src:src:math:matrix.rs]{generado/fuente/src/math/matrix.rs}
\lstinputlisting[language=Rust,style=anexo,caption={Archivo \texttt{src/math/rng.rs}},label=cod:src:src:math:rng.rs]{generado/fuente/src/math/rng.rs}
\section{Módulo \texttt{nn}: red neuronal}
Esta sección contiene \texttt{src/nn/mod.rs} (19 líneas), \texttt{src/nn/activation.rs} (136 líneas), \texttt{src/nn/loss.rs} (45 líneas), \texttt{src/nn/neuron.rs} (55 líneas), \texttt{src/nn/layer.rs} (205 líneas), \texttt{src/nn/network.rs} (234 líneas).
\lstinputlisting[language=Rust,style=anexo,caption={Archivo \texttt{src/nn/mod.rs}},label=cod:src:src:nn:mod.rs]{generado/fuente/src/nn/mod.rs}
\lstinputlisting[language=Rust,style=anexo,caption={Archivo \texttt{src/nn/activation.rs}},label=cod:src:src:nn:activation.rs]{generado/fuente/src/nn/activation.rs}
\lstinputlisting[language=Rust,style=anexo,caption={Archivo \texttt{src/nn/loss.rs}},label=cod:src:src:nn:loss.rs]{generado/fuente/src/nn/loss.rs}
\lstinputlisting[language=Rust,style=anexo,caption={Archivo \texttt{src/nn/neuron.rs}},label=cod:src:src:nn:neuron.rs]{generado/fuente/src/nn/neuron.rs}
\lstinputlisting[language=Rust,style=anexo,caption={Archivo \texttt{src/nn/layer.rs}},label=cod:src:src:nn:layer.rs]{generado/fuente/src/nn/layer.rs}
\lstinputlisting[language=Rust,style=anexo,caption={Archivo \texttt{src/nn/network.rs}},label=cod:src:src:nn:network.rs]{generado/fuente/src/nn/network.rs}
}

\chapter{Inventario de diagramas editables}\label{anx:drawio}
Los diagramas del documento se elaboraron en formato \texttt{.drawio} (diagrams.net), editable y abierto, y se exportaron a PNG a 300 DPI con el script \path{diagramas/generar_diagramas.py}. La Tabla~\ref{tab:drawio} presenta el inventario completo.

\begin{table}[!htbp]
\caption{Inventario de archivos \texttt{.drawio} del proyecto}\label{tab:drawio}
\small
\begin{tabularx}{\textwidth}{@{}>{\raggedright}p{5.4cm} >{\raggedright}X >{\raggedright\arraybackslash}p{1.6cm}@{}}
\toprule
Archivo fuente y exportación & Contenido & Figura \\
\midrule
\path{01_arbol_problemas.drawio} \newline \path{diag01_arbol_problemas.png} & Causas, problema central y efectos & \ref{fig:arbol} \\
\path{02_arquitectura_modulos_rust.drawio} \newline \path{diag02_arquitectura_rust.png} & Módulos \texttt{math}, \texttt{nn}, \texttt{experiments} y CLI & \ref{fig:arquitectura} \\
\path{03_flujo_forward_backprop.drawio} \newline \path{diag03_flujo_entrenamiento.png} & Algoritmo de entrenamiento & \ref{fig:flujo} \\
\path{04_memoria_ownership_rust.drawio} \newline \path{diag04_memoria_rust.png} & Pila, montón, propiedad y préstamos & \ref{fig:memoria} \\
\path{05_fases_metodologia.drawio} \newline \path{diag05_fases_metodologia.png} & Fases y etapas de la metodología & \ref{fig:fases} \\
\bottomrule
\end{tabularx}
\fuente{Archivos fuente en \path{diagramas/}; exportaciones en \path{tesis/figuras/}. Pueden abrirse en \url{https://app.diagrams.net} o en la aplicación de escritorio draw.io.}
\end{table}

Las once gráficas de resultados se generan con \path{scripts/generar_graficas.py} a partir de los CSV, también a 300 DPI, y se conservan en \path{tesis/figuras/} con prefijo \texttt{fig}. Al modificar un diagrama o volver a ejecutar los experimentos, basta repetir los comandos del Código~\ref{cod:regenerar} para actualizar todas las figuras de este documento.

\chapter{Integrantes, ORCID y roles}\label{anx:roles}
El equipo de investigación está conformado por cinco estudiantes de la FICCT inscritos en la asignatura \Materia{} (\Sigla), grupo \Grupo. La Tabla~\ref{tab:roles} presenta sus datos de identificación y su rol en la Feria Facultativa.

\begin{table}[!htbp]
\caption{Integrantes del equipo, identificadores ORCID y roles}\label{tab:roles}
\small\setlength{\tabcolsep}{4pt}
\begin{tabularx}{\textwidth}{@{}>{\raggedright}X l >{\raggedright}p{2.4cm} l >{\raggedright\arraybackslash}p{2.6cm}@{}}
\toprule
Nombre completo & Registro & Carrera & ORCID & Rol \\
\midrule
Ronal Miranda Vasquez & 226159787 & Ing. Informática & 0009-0003-9735-4920 & Investigador -- Expositor \\
Oscar Alejandro Nava Orihuela & 226037215 & Ing. Informática & 0009-0001-0485-3030 & Investigador -- Expositor \\
Carla Olivia Cano Quispe & 224026879 & Ing. Informática & 0009-0009-5668-8290 & Investigadora -- Expositora \\
Diana Nayely Diaz Mamani & 225028190 & Ing. Sistemas & 0009-0003-0550-4996 & Investigadora -- Expositora \\
Julio Cesar Piuca Achocalla & 224028782 & Ing. Informática & 0009-0004-0056-5201 & Investigador -- Expositor \\
\bottomrule
\end{tabularx}
\fuente{Docente tutor: \DocenteTutor. Coordinador del IICCT: \CoordinadorIICCT.}
\end{table}

Las responsabilidades del proyecto se distribuyeron por áreas de trabajo en las que participó todo el equipo: implementación del motor en Rust por fases, diseño y ejecución de las pruebas unitarias, ejecución y análisis de los experimentos, elaboración de gráficas y diagramas, redacción del documento y preparación del material y de la exposición en la feria. Cada integrante se encuentra en condiciones de explicar y demostrar cualquiera de los componentes del sistema durante la exposición.

\chapter{Carta de Compromiso}\label{anx:carta}
\begin{center}
\includegraphics[width=0.85\textwidth,height=0.78\textheight,keepaspectratio]{\CartaCompromiso}
\end{center}
